\chapter{The Colour of Gold}
\label{ch:colour}

Chapter~\ref{ch:intro} promised that the program would explain why gold is
yellow. This chapter keeps the promise as far as an honest calculation
can: it derives, implements and tests the complete chain
\[
\text{bands and wavefunctions}
\;\to\; \varepsilon(\omega)
\;\to\; R(\omega)
\;\to\; \text{colour seen by a human eye},
\]
and then looks critically at what comes out. Every step is a
post-processing of the fixed-potential calculation of
Chapter~\ref{ch:conductivity}: no new self-consistent calculation is
needed.

\section{What ``colour'' means physically}

A thick piece of metal transmits nothing; what you see is the light it
reflects. Its colour is therefore determined by the
\textbf{reflectivity} $R(\omega)$, the fraction of the incident intensity
that is reflected at each photon energy $\hbar\omega$. Visible light spans
$\hbar\omega = 1.6$\,eV (red, $\lambda = 780$\,nm) to $3.3$\,eV (violet,
$380$\,nm); the conversion is $\lambda[\mathrm{nm}] =
1239.84/\hbar\omega[\mathrm{eV}]$. A metal that reflects all visible
frequencies equally is white-grey (silver, aluminium). Gold reflects red,
orange and yellow light almost completely and blue light poorly; the eye
reads that mixture as ``golden yellow''. Copper cuts off at a lower
energy still and looks reddish.

The reflectivity in turn is fixed by a single complex function, the
\textbf{dielectric function} $\varepsilon(\omega) = \varepsilon_1(\omega)
+ \ii\varepsilon_2(\omega)$, which describes how the electrons respond to
an oscillating electric field. So the task is to compute
$\varepsilon(\omega)$.

\section{The dielectric function from the band structure}

An electromagnetic wave of frequency $\omega$ can do two things to the
electrons of a metal.

\paragraph{Intraband response: the Drude term.} It can accelerate the
electrons at the Fermi surface, exactly as a static field does in
Chapter~\ref{ch:conductivity}. Repeating that derivation with
$\mathbf E\propto\ee^{-\ii\omega t}$ gives the conductivity
$\sigma(\omega) = \sigma(0)/(1-\ii\omega\tau)$, and since a conductivity
contributes $4\pi\ii\sigma/\omega$ to the dielectric function,
\begin{equation}
\varepsilon_{\mathrm{intra}}(\omega) = -\frac{\omega_p^2}{\omega\,(\omega +
\ii\gamma)},
\qquad \gamma = 1/\tau ,
\label{eq:drude}
\end{equation}
with the plasma frequency of Eq.~\eqref{eq:plasma}. For
$\gamma\ll\omega\ll\omega_p$ this is large and negative: the free
electrons move so as to cancel the field inside the metal, and the wave is
reflected. This is why metals are shiny.

\paragraph{Interband response: absorption.} It can lift an electron from
an occupied state $n\kk$ to an empty state $m\kk$ of another band. Light
carries negligible momentum on the scale of the Brillouin zone, so $\kk$
does not change: transitions are \emph{vertical} in a band-structure plot.
The photon is absorbed if its energy matches,
$\hbar\omega = \varepsilon_{m\kk}-\varepsilon_{n\kk}\equiv\omega_{mn}$,
with a rate given by Fermi's golden rule. The perturbation is the coupling
of the field to the electron's velocity, so the rate is proportional to
$|\mathbf v_{mn}|^2$, the off-diagonal velocity matrix elements of
Eq.~\eqref{eq:vmn}. Summing the absorbed power over all transitions and
dividing by the energy flux of the wave gives the absorptive part of the
dielectric function (cubic crystal, atomic units, one electron per spinor
state):
\begin{equation}
\boxed{\;
\varepsilon_2^{\mathrm{inter}}(\omega) = \frac{4\pi^2}{\Omega}
\sum_{\kk} w_\kk \sum_{n,m}\big(f_{n\kk}-f_{m\kk}\big)\,
\frac{|\mathbf v_{mn}(\kk)|^2}{3\,\omega_{mn}^2}\;
\delta\big(\omega - \omega_{mn}\big),
\qquad \omega_{mn} > 0 .\;}
\label{eq:eps2}
\end{equation}
Read it as a recipe: for every $\kk$ and every pair (occupied $n$, empty
$m$), put a spike at the transition energy $\omega_{mn}$ with strength
$|\mathbf v_{mn}|^2/3\omega_{mn}^2$. The factor $f_n - f_m$ enforces the
Pauli principle (from a filled to an empty state); $\tfrac13|\mathbf
v|^2$ is the average over the three polarization directions, which are
equivalent in a cubic crystal.

\paragraph{The real part.} Causality ties the real part of a response
function to its imaginary part (the Kramers--Kronig relation,
$\varepsilon_1(\omega) = 1 + \tfrac{2}{\pi}\,\mathcal P\!\!\int_0^\infty
\frac{\omega'\varepsilon_2(\omega')}{\omega'^2-\omega^2}\,\dd\omega'$).
Inserting Eq.~\eqref{eq:eps2} does the integral analytically and gives
both parts in one complex formula, which is what you should program:
\begin{equation}
\varepsilon_{\mathrm{inter}}(\omega) = 1 + \frac{8\pi}{\Omega}
\sum_\kk w_\kk \sum_{\substack{n,m\\ \omega_{mn}>0}}
\big(f_{n\kk}-f_{m\kk}\big)\,
\frac{|\mathbf v_{mn}(\kk)|^2/3}
{\omega_{mn}\,\big[\omega_{mn}^2 - (\omega + \ii\eta)^2\big]} .
\label{eq:epsinter}
\end{equation}
The small positive $\eta$ turns each $\delta$-spike of
Eq.~\eqref{eq:eps2} into a Lorentzian of width $\eta$ (check: for
$\eta\to0$, $\mathrm{Im}\,[\omega_{mn}^2-(\omega+\ii\eta)^2]^{-1} \to
\pi\,\delta(\omega-\omega_{mn})/2\omega_{mn}$, which reproduces
Eq.~\eqref{eq:eps2}). It represents the finite lifetime of the excited
states and, in practice, smooths the result of a finite $\kk$-mesh. The
total dielectric function is
\begin{equation}
\varepsilon(\omega) = \varepsilon_{\mathrm{inter}}(\omega)
+ \varepsilon_{\mathrm{intra}}(\omega) .
\label{eq:epstot}
\end{equation}

\paragraph{What has been approximated.} Equations
\eqref{eq:drude}--\eqref{eq:epstot} are the \emph{independent-particle
approximation}: (i) Kohn--Sham eigenvalue differences are used as
excitation energies --- an assumption, not a theorem (see the
interpretation warnings of Chapter~\ref{ch:hkks}); (ii) the electron and
the hole it leaves behind are assumed not to interact, which is a good
approximation in a metal, where they are screened; (iii) the microscopic
fields that vary on the scale of the unit cell (``local fields'') are
neglected; (iv) only as many empty bands as were computed are included,
which slightly underestimates $\varepsilon_1$. For the colour of a noble
metal (i) is by far the most important, as we shall see.

\section{From the dielectric function to the reflectivity}

Maxwell's equations for a wave hitting a flat surface at normal incidence
give the reflected fraction in terms of the complex refractive index
$\tilde n = \sqrt{\varepsilon}$:
\begin{equation}
R(\omega) = \left|\frac{\sqrt{\varepsilon(\omega)} - 1}
{\sqrt{\varepsilon(\omega)} + 1}\right|^2 .
\label{eq:reflectivity}
\end{equation}
Two limits explain everything that follows. If $\varepsilon$ is real and
negative, $\sqrt\varepsilon$ is purely imaginary and $R = 1$: total
reflection, the Drude regime. If absorption sets in
($\varepsilon_2$ no longer small compared with $|\varepsilon_1|$), $R$
drops. \emph{The colour of a noble metal is decided by the photon energy
at which interband absorption begins}: below that threshold the metal is
a near-perfect mirror, above it a mediocre one. In gold the threshold is
the energy needed to lift an electron from the top of the $5d$ bands to
the Fermi level (Fig.~\ref{fig:bands}). Experimentally it lies at about
$2.4$\,eV ($\approx 520$\,nm, blue-green), so blue light is absorbed (the measured reflectivity falls through $0.7$ at
$2.33$\,eV for gold, $2.16$\,eV for copper and $3.71$\,eV for silver,
interpolating between the tabulated points of Johnson and Christy). In silver, whose $4d$ bands lie deeper, it
is at about $3.8$\,eV, in the ultraviolet --- so silver reflects the whole
visible spectrum and is white. The difference between the two is
relativity (Chapter~\ref{ch:soc}): the relativistic contraction of the
$6s$ shell and expansion of the $5d$ shell bring gold's $d$ bands close
to the Fermi level.

\section{From the reflectivity to a colour}

Human colour vision has three kinds of receptors. Any spectrum of light
$S(\lambda)$ is perceived through just three numbers, the
\emph{tristimulus values}
\begin{equation}
X = \int S(\lambda)\,\bar x(\lambda)\,\dd\lambda,\qquad
Y = \int S(\lambda)\,\bar y(\lambda)\,\dd\lambda,\qquad
Z = \int S(\lambda)\,\bar z(\lambda)\,\dd\lambda ,
\end{equation}
where $\bar x,\bar y,\bar z$ are the standardized \emph{colour-matching
functions} of the CIE 1931 standard observer ($Y$ measures brightness).
For a reflecting object under an illuminant of spectrum $I(\lambda)$, the
light reaching the eye is $S = I\,R$. Four ingredients make this
computable in a few lines.

\paragraph{1. Colour-matching functions.} They are tabulated, but an
accurate analytic fit exists (Wyman, Sloan and Shirley, 2013). With
$\lambda$ in nm and the asymmetric Gaussian
$g(\lambda;\mu,\sigma_1,\sigma_2) =
\exp\!\big[-\tfrac12(\lambda-\mu)^2/\sigma^2\big]$, $\sigma = \sigma_1$
for $\lambda<\mu$ and $\sigma_2$ otherwise:
\begin{equation}
\begin{aligned}
\bar x(\lambda) &= 1.056\,g(\lambda;599.8,37.9,31.0)
+ 0.362\,g(\lambda;442.0,16.0,26.7) - 0.065\,g(\lambda;501.1,20.4,26.2),\\
\bar y(\lambda) &= 0.821\,g(\lambda;568.8,46.9,40.5)
+ 0.286\,g(\lambda;530.9,16.3,31.1),\\
\bar z(\lambda) &= 1.217\,g(\lambda;437.0,11.8,36.0)
+ 0.681\,g(\lambda;459.0,26.0,13.8).
\end{aligned}
\label{eq:cmf}
\end{equation}
(Checked against the CIE table: the largest deviations are $0.015$,
$0.007$ and $0.024$ for functions that peak at $1.06$, $1.00$ and
$1.78$.)

\paragraph{2. Illuminant.} Standard daylight (CIE D65) is approximated here
by a black body at $T = 6500$\,K (the true D65 spectrum differs, mostly
in the violet; tabulated D65 may be substituted, but all
reference colours in this document assume the black body), $I(\lambda)\propto\lambda^{-5}/
[\exp(c_2/\lambda T)-1]$ with $c_2 = 1.438777\times10^{7}$\,nm\,K.

\paragraph{3. Normalization.} Divide by the values a perfect mirror
($R = 1$) would give, so that white stays white whatever the illuminant,
and express the result relative to the white point of the display
standard:
\begin{equation}
X = 0.95047\,\frac{\int I R\,\bar x\,\dd\lambda}{\int I\,\bar x\,\dd\lambda},
\qquad
Y = \frac{\int I R\,\bar y\,\dd\lambda}{\int I\,\bar y\,\dd\lambda},
\qquad
Z = 1.08883\,\frac{\int I R\,\bar z\,\dd\lambda}{\int I\,\bar z\,\dd\lambda}.
\label{eq:xyz}
\end{equation}
The integrals run from 380 to 780\,nm; a sum over 1\,nm steps suffices.
Scaling each channel separately in this way is the simplest (``XYZ
scaling'') form of \emph{chromatic adaptation}; colour scientists use
more refined transforms (e.g.\ Bradford), which shift the final sRGB
values by a few units but change nothing in the physics.

\paragraph{4. Display colour (sRGB).} Screens use the sRGB standard. The
linear red, green and blue intensities are
\begin{equation}
\begin{pmatrix} r\\ g\\ b\end{pmatrix} =
\begin{pmatrix}
\phantom{-}3.2406 & -1.5372 & -0.4986\\
-0.9689 & \phantom{-}1.8758 & \phantom{-}0.0415\\
\phantom{-}0.0557 & -0.2040 & \phantom{-}1.0570
\end{pmatrix}
\begin{pmatrix} X\\ Y\\ Z\end{pmatrix},
\label{eq:srgbmatrix}
\end{equation}
each clipped to $[0,1]$ and then encoded with the sRGB transfer function
$c \mapsto 12.92\,c$ for $c\le0.0031308$ and $1.055\,c^{1/2.4}-0.055$
otherwise; multiply by 255 and round to obtain the familiar 8-bit triple.
Two unit tests: $R\equiv1$ must give $(255,255,255)$, and $R\equiv0.5$
must give $(188,188,188)$.

\begin{algorithm}[H]
\caption{Dielectric function, reflectivity and colour}
\label{alg:colour}
\begin{algorithmic}[1]
\State \textbf{input:} $\varepsilon_{n\kk}$, $w_\kk$,
$V_{mn}(\kk) = \tfrac13\sum_\alpha|v^\alpha_{mn}|^2$ and $\omega_p$ from
Algorithm~\ref{alg:drude};\ $\mu$;\ $\eta$, $\gamma$
\State $f_{n\kk}\gets 1$ if $\varepsilon_{n\kk}<\mu$ else $0$
\Comment{sharp occupations: no smearing here}
\State collect all transitions with $\omega_{mn} =
\varepsilon_{m\kk}-\varepsilon_{n\kk} > 10^{-5}$\,Ha and $f_n - f_m > 0$:
energies $\omega_t$, strengths $S_t = w_\kk\,(f_n-f_m)\,V_{mn}(\kk)$
\For{each $\omega$ on a grid, $0 < \hbar\omega \le 8$\,eV, step $0.01$\,eV}
\State $\varepsilon(\omega) \gets 1 + \dfrac{8\pi}{\Omega}\sum_t
\dfrac{S_t}{\omega_t\,[\omega_t^2-(\omega+\ii\eta)^2]}
\;-\;\dfrac{\omega_p^2}{\omega(\omega+\ii\gamma)}$
\State $R(\omega)\gets\big|(\sqrt{\varepsilon}-1)/(\sqrt{\varepsilon}+1)\big|^2$
\EndFor
\State interpolate $R$ to $\lambda = 380,381,\dots,780$\,nm
($\hbar\omega = 1239.84/\lambda$ eV)
\State $X,Y,Z$ from Eqs.~\eqref{eq:cmf}--\eqref{eq:xyz};\quad sRGB from
Eq.~\eqref{eq:srgbmatrix} and the transfer function
\end{algorithmic}
\end{algorithm}

The transitions within a Kramers pair have $\omega_{mn} = 0$ and are
excluded by the threshold in line 3; they are the intraband term. The
sum over $\kk$ uses the irreducible wedge with its weights, as in
Chapter~\ref{ch:conductivity}.

\section{Results: the computed spectrum of gold}
\label{sec:colourresults}

The reference calculation uses the fixed self-consistent potential of
Case~B (LDA, spin--orbit coupling on), the irreducible wedge of a
$\Gamma$-centered $40^3$ mesh (\NKFINE{} points, 22 spinor bands, about
\NTRANS\,000 transitions), $\eta = 0.1$\,eV and $\hbar\gamma = 0.07$\,eV
(a typical value from optical measurements on gold; the colour is
insensitive to it). Figure~\ref{fig:optics} shows the result.

\begin{figure}[tbp]
\centering
\pgfplotsset{optplot/.style={width=0.5\textwidth, height=0.42\textwidth,
  xmin=0.5, xmax=6, xlabel={$\hbar\omega$ (eV)},
  tick label style={font=\small}, label style={font=\small},
  title style={font=\small}, legend style={font=\scriptsize, draw=none, fill=none}, legend cell align={left},
  no markers}}
\begin{tikzpicture}
\begin{axis}[optplot, ymin=0, ymax=13, title={Interband absorption $\varepsilon_2^{\mathrm{inter}}$},
  legend pos=north east]
\fill[yellow!25] (axis cs:1.6,0) rectangle (axis cs:3.3,13);
\addplot[red!70!black, thick] table[x index=0, y index=3] {figdata/optics_spec.dat};
\addplot[black!55, thin] table[x index=0, y index=9] {figdata/optics_spec.dat};
\addplot[black, only marks, mark=o, mark size=1.2pt] table[x index=0, y index=2] {figdata/optics_exp.dat};
\legend{$40^3$ mesh, $24^3$ mesh, measured $\varepsilon_2$}
\end{axis}
\end{tikzpicture}\hfill
\begin{tikzpicture}
\begin{axis}[optplot, ymin=0, ymax=1.05, title={Reflectivity $R$},
  legend style={at={(0.97,0.86)}, anchor=north east}]
\fill[yellow!25] (axis cs:1.6,0) rectangle (axis cs:3.3,1.05);
\addplot[red!70!black, thick] table[x index=0, y index=4] {figdata/optics_spec.dat};
\addplot[blue!60!black, thick, dashed] table[x index=0, y index=6] {figdata/optics_spec.dat};
\addplot[black!55, thin] table[x index=0, y index=7] {figdata/optics_spec.dat};
\addplot[black, only marks, mark=o, mark size=1.2pt] table[x index=0, y index=3] {figdata/optics_exp.dat};
\legend{computed (LDA), transitions $+0.75$\,eV, Drude term only, measured}
\end{axis}
\end{tikzpicture}
\caption{Optical response of gold from the reference calculation. The
shaded band is the visible range ($1.6$--$3.3$\,eV). Left: the interband
part of $\varepsilon_2$, Eq.~\eqref{eq:eps2}, for two $\kk$-meshes; the
remaining differences are the fine structure a finite mesh cannot
resolve. Right: reflectivity, Eq.~\eqref{eq:reflectivity}. Without
interband transitions gold would be a featureless mirror (grey line);
the $5d\to6s$ transitions cut the reflectivity down inside the visible
range (red). Dashed: the same calculation with all interband transition
energies raised by $0.75$\,eV, to show what a correctly placed absorption
edge does (see text). Circles: measured optical constants of Johnson and
Christy (1972); the measured $\varepsilon_2$ in the left panel is the
total, including the Drude part, which dominates below about $1.3$\,eV
(the measured point at $0.64$\,eV, $\varepsilon_2 = 25$, lies above the
frame).}
\label{fig:optics}
\end{figure}

\paragraph{What the calculation gets right.} Below about $1.3$\,eV
interband absorption is weak and gold is an excellent mirror
($R > 0.95$): this is the Drude regime, governed by $\omega_p$ alone.
Then the $5d\to6s$ transitions set in. $\varepsilon_2^{\mathrm{inter}}$
rises steeply, and the reflectivity collapses from above $0.9$ to below
$0.5$ within half an electron volt. Above the edge it stays between
$0.4$ and $0.5$ throughout the visible, close to the measured reflectivity
of gold for blue light (\REXPb{} at 450\,nm). The mechanism of gold's colour ---
a reflectivity edge in the visible, created by the $d$ bands lying only a
couple of eV below the Fermi level --- comes out of the band structure
with no adjustable parameter. So does the integrated strength: the
interband transitions below 8\,eV account for \NEFFINTER{} electrons per
atom in the sum rule, on top of the \NEFF{} conduction electrons.

\paragraph{What it gets wrong, and why.} The \emph{position} of the edge.
The computed reflectivity falls through $0.7$ at \EDGE\,eV; in real gold
(circles in Fig.~\ref{fig:optics}) the corresponding drop is at
$2.33$\,eV. The cause is not the
optics but the band structure it is fed: the $5d$ bands of this
calculation lie too close to the Fermi level. Their upper edge is at
$-1.05$\,eV at $X$ and $-1.63$\,eV at $L$ (Table~\ref{tab:goldlevels}),
whereas experiment places them $0.5$--$0.9$\,eV deeper
(Table~\ref{tab:dbandexp}). This is the ``known shortcoming of LDA
and GGA for the $d$ bands of noble metals'' met in
Chapter~\ref{ch:soc}, made somewhat worse here by the 11-electron
pseudopotential, and it is approximation (i) of the list above at work:
Kohn--Sham eigenvalues are not excitation energies.

The consequence for the colour is drastic, because the edge sits where the
eye is most sensitive. With the edge at \EDGE\,eV instead of $2.33$\,eV,
green and yellow light is absorbed as strongly as blue, and only red is
reflected well (Table~\ref{tab:colours}).

\begin{table}[tbp]
\centering
\begin{tabular}{@{}lcccc@{}}
\toprule
 & \multicolumn{3}{c}{reflectivity at} & colour\\
 & 450\,nm & 550\,nm & 650\,nm & sRGB\\
\midrule
Drude term only (no $d$ bands) & \RDRb & \RDRg & \RDRr &
\fcolorbox{black}{goldDR}{\rule{0pt}{1.6ex}\hspace{2.2em}}\ \ $(\RGBDR)$\\
computed, LDA & \RLDAb & \RLDAg & \RLDAr &
\fcolorbox{black}{goldLDA}{\rule{0pt}{1.6ex}\hspace{2.2em}}\ \ $(\RGBLDA)$\\
interband transitions $+0.5$\,eV & \RSHab & \RSHag & \RSHar &
\fcolorbox{black}{goldSHa}{\rule{0pt}{1.6ex}\hspace{2.2em}}\ \ $(\RGBSHa)$\\
interband transitions $+0.75$\,eV & \RSHbb & \RSHbg & \RSHbr &
\fcolorbox{black}{goldSHb}{\rule{0pt}{1.6ex}\hspace{2.2em}}\ \ $(\RGBSHb)$\\
measured (Johnson and Christy) & \REXPb & \REXPg & \REXPr &
\fcolorbox{black}{goldEXP}{\rule{0pt}{1.6ex}\hspace{2.2em}}\ \ $(\RGBEXP)$\\
\bottomrule
\end{tabular}
\caption{Reflectivity at three wavelengths and the resulting display colour
(swatch and 8-bit sRGB triple) for the reference calculation, two
illustrative shifts of the interband transition energies, the Drude term
alone, and the measured optical constants.}
\label{tab:colours}
\end{table}

The last row of Table~\ref{tab:colours} is itself a test of the colour chain of
Algorithm~\ref{alg:colour}: feeding \emph{measured} reflectivities
through it gives
\definecolor{agEXP}{RGB}{254,253,252}\definecolor{cuEXP}{RGB}{246,208,191}%
\fcolorbox{black}{goldEXP}{\rule{0pt}{1.6ex}\hspace{1.4em}} $(\RGBEXP)$
for gold,
\fcolorbox{black}{agEXP}{\rule{0pt}{1.6ex}\hspace{1.4em}} $(254,253,252)$
for silver and
\fcolorbox{black}{cuEXP}{\rule{0pt}{1.6ex}\hspace{1.4em}} $(246,208,191)$
for copper (optical constants of Johnson and Christy for all three) ---
the colours of the three metals. The parameter-free LDA result, by
contrast, is a pale, pinkish metal --- warmer than
silver, far from yellow: the calculation predicts \emph{that} gold is
coloured and \emph{why}, but not the right hue. The two ``shifted'' rows
are an illustration, not a prediction: every interband transition energy
$\omega_{mn}$ in the denominators of Eq.~\eqref{eq:epsinter} is raised by
a constant $\Delta$ while the strengths $|\mathbf v_{mn}|^2/\omega_{mn}^2$
are kept (in Algorithm~\ref{alg:colour}, replace
$S_t/(\omega_t[\omega_t^2-\dots])$ by
$(S_t/\omega_t^2)\,(\omega_t+\Delta)/[(\omega_t+\Delta)^2-\dots]$). As
the edge moves toward its experimental position, green and yellow light
is reflected again and the colour turns from pink through salmon to an
orange gold. A rigid shift is a caricature of the true correction (it
moves the edge by less than $\Delta$, because $\varepsilon_1$ changes
with it, and it leaves the shape of the spectrum unchanged), but it makes
the essential point: \emph{the colour of gold is a sensitive measurement
of the energy of its $d$ bands}, more sensitive than the accuracy of
LDA.

\begin{table}[H]
\centering\small
\begin{tabular}{@{}lccc@{}}
\toprule
Level (Kramers pair) & this calculation & PBE, semicore & experiment\\
\midrule
$\Gamma$: $d$ levels & $-5.45$, $-4.24$, $-3.02$ & $-5.67$, $-4.46$, $-3.27$ & $-6.0$, $-4.6$, $-3.7$\\
$L$: upper $d$ levels (pairs 4, 5) & $-2.36$, $-1.63$ & $-2.61$, $-1.90$ & $-3.2$, $-2.3$\\
$L$: $p$-like level (pair 6) & $-1.22$ & $-1.37$ & $-1.0$\\
$L$: $s$-like level (pair 7) & $+3.13$ & $+2.93$ & $+3.6$\\
\bottomrule
\end{tabular}
\caption{Band energies of gold in eV relative to $E_F$: the reference
calculation of this document (LDA, 11-electron pseudopotential), a
published all-electron-quality PBE calculation with semicore states, and
angle-resolved photoemission and optical experiments (the last two columns
from Rangel \emph{et al.}, Phys.\ Rev.\ B \textbf{86}, 125125 (2012),
Table~II; experimental values from different measurements scatter by
about $0.1$\,eV). The $d$ levels of both density-functional calculations
are too shallow; those of this calculation are a further
$0.2$--$0.3$\,eV shallower than the published PBE ones.}
\label{tab:dbandexp}
\end{table}

The \emph{strength} of the absorption is also overestimated: at 3\,eV the
computed $\varepsilon_2$ is \EPSCALC{} against a measured \EPSEXP. This
too is typical of the independent-particle approximation with LDA bands
for the noble metals.

\paragraph{How one would do better.} In order of increasing effort:
(a) the 19-electron pseudopotential of Appendix~\ref{app:q19} and the PBE
functional, which change the $d$-band position by tenths of an eV;
(b) a functional or method that describes localized $d$ states better
than LDA and GGA --- hybrid functionals or DFT$+U$; (c) many-body
perturbation theory ($GW$), in which quasiparticle energies replace
Kohn--Sham eigenvalues and which, for gold, widens the $5d$--$6sp$
separation by $0.35$--$0.75$\,eV (Rangel \emph{et al.}, cited above). All three change the
\emph{input} to this chapter; Algorithm~\ref{alg:colour} stays as it is.
None of them has been carried out for this document.

\paragraph{The role of relativity.} The argument of
Chapter~\ref{ch:soc} can now be made quantitative in the opposite
direction. Everything that distinguishes the red curve of
Fig.~\ref{fig:optics} from the grey one is the $5d$ band lying within
2\,eV of the Fermi level, and that proximity is the relativistic
contraction of the $6s$ shell and expansion of the $5d$ shell, built into
the pseudopotential. Silver, with the same crystal structure and nearly
the same lattice constant but weaker relativistic effects, has its $4d$
bands about 2\,eV deeper: its edge lies in the ultraviolet and its
visible reflectivity is the grey curve. Spin--orbit coupling itself plays
a secondary role for the colour: it moves the upper $d$-band edge by a few
tenths of an eV (compare the two halves of Table~\ref{tab:goldlevels}),
and the reflectivity edge with it (\EDGENOSOC\,eV without SOC against
\EDGESOC\,eV with SOC on the same $16^3$ mesh).

\paragraph{Numerical convergence.} The sRGB values change by up to about
ten units out of 255 when $\eta$ is varied between $0.05$ and $0.2$\,eV
($(\RGBETA)$ for $\eta = 0.2$\,eV) and by up to about fifteen units
between the $24^3$ and $40^3$ meshes ($(\RGBCO)$ on the coarser one); the Drude
damping $\gamma$ has no visible effect. These numerical uncertainties are
real but small compared with the effect of the misplaced $d$ bands.

\section{Sanity checks}

\begin{enumerate}[leftmargin=2em]
\item \textbf{Velocity matrix elements}: the Hellmann--Feynman test of
Section~\ref{sec:velocity}.
\item \textbf{Drude limit}: with the interband sum switched off,
$R(\omega)$ must be within a few percent of 1 for
$\gamma\ll\omega<\omega_p$ and drop steeply once $\omega$ exceeds
$\omega_p$; the colour must come out as a neutral near-white
(the reference calculation gives sRGB $(\RGBDR)$).
\item \textbf{Sum rule}: $N_{\mathrm{eff}}(\omega_{\max}) =
\frac{\Omega}{2\pi^2}\int_0^{\omega_{\max}}\omega\,\varepsilon_2(\omega)\,
\dd\omega$ counts the electrons per cell that have absorbed below
$\omega_{\max}$. The intraband term alone must give
$\omega_p^2\Omega/4\pi$ (about one electron), and the total must stay
below the number of valence electrons, 11.
\item \textbf{Colour chain}: white and 50\% grey as above; a reflectivity
that is 1 above 600\,nm and 0 below must come out red, and so on.
\item \textbf{Kramers--Kronig}: $\varepsilon_1$ from
Eq.~\eqref{eq:epsinter} must agree with a numerical Kramers--Kronig
transform of $\varepsilon_2$.
\item \textbf{Convergence}: repeat with a denser $\kk$-mesh and with
$\eta$ halved. The positions of the absorption edge and of the
reflectivity drop must not move; fine structure will.
\end{enumerate}

\begin{gpubox}{GPU/JAX notes}
The expensive part --- bands and velocity matrices on the dense mesh ---
was done in Chapter~\ref{ch:conductivity}. Everything in this chapter is a
sum over a list of a few hundred thousand transitions for each of a few
hundred frequencies: an outer product that fits one \texttt{vmap} over $\omega$,
or a few seconds of NumPy on the host. A faster and equally accurate
route is to histogram the strengths $S_t/\omega_t^2$ on a fine energy
grid once and convolve.
\end{gpubox}
